\section{Hasil Pengujian}
\TemplateTodo{Sajikan hasil utama sesuai skenario pada Bab III. Tabel harus memiliki judul jelas, satuan, dan keterangan yang cukup.}

Tabel~\ref{tab:confusion} menunjukkan \emph{confusion matrix} dari pengujian menggunakan 10-fold cross validation pada metode Na\"ive Bayes Classifier.

\begin{table}[H]
\centering
\caption{Confusion matrix hasil pengujian NBC}
\label{tab:confusion}
\begin{tabular}{@{}lcc@{}}
\toprule
 & \multicolumn{2}{c}{Prediksi} \\
\cmidrule{2-3}
Aktual & Positif & Negatif \\
\midrule
Positif & 2.132 (TP) & 287 (FN) \\
Negatif & 398 (FP) & 2.183 (TN) \\
\bottomrule
\end{tabular}
\end{table}

Tabel~\ref{tab:comparison} menyajikan perbandingan metrik evaluasi antara Na\"ive Bayes Classifier dan Support Vector Machine.

\begin{table}[H]
\centering
\caption{Perbandingan performa NBC dan SVM}
\label{tab:comparison}
\begin{tabular}{@{}lcc@{}}
\toprule
Metrik & NBC & SVM \\
\midrule
Akurasi & 86,3\% & 82,1\% \\
Precision & 84,7\% & 80,5\% \\
Recall & 88,1\% & 83,9\% \\
F1-Score & 86,4\% & 82,2\% \\
\bottomrule
\end{tabular}
\end{table}
